\documentclass[10pt,a4paper]{amsart}
\usepackage[margin=19mm]{geometry}
\usepackage{amsmath,amssymb,mathtools}
\usepackage[scr=rsfs,cal=euler]{mathalfa}
\usepackage{xfrac}
\usepackage[unicode,hidelinks,bookmarks=false]{hyperref}
\newcommand{\ie}{i.e.\ }
\newcommand{\cf}{cf.\ }
\newcommand{\resp}{respectively\ }
\newcommand{\Subsection}{\subsection}
\providecommand{\nobreakdash}{\nobreak\hbox{-}}
\setlength{\emergencystretch}{2em}
\multlinegap0pt
\usepackage{amssymb}
\usepackage{float}
\newtheorem{lemma}{Lemma}
\title{Real ideal points, Conway spheres, and left-orderable Dehn fillings}
\author{Yi Wang}
\date{Source: arXiv:2608.26564v1 (CC BY 4.0)}
\begin{document}
\maketitle

\noindent\emph{Source and licence.} Real ideal points, Conway spheres, and left-orderable Dehn fillings. Version arXiv:2608.26564v1 (CC BY 4.0), distributed by arXiv under the Creative Commons Attribution 4.0 International licence, http://creativecommons.org/licenses/by/4.0/, as recorded in the arXiv licence metadata for this exact version and shown on the abstract page https://arxiv.org/abs/2608.26564v1. Full licence notice: \texttt{licenses/arxiv-2608.26564-CC-BY-4.0.txt}.\par\smallskip

\noindent\textit{Editorial scope.} Counted unit: the article's lemma lma:commutatorasymptotics (source0.tex lines 2388--2404). The article's complete original proof of it is delivered with it (source0.tex lines 2406--2475, one \texttt{proof} environment of 70 lines). No mathematics is written, repaired or re-derived: the statement and the proof are the article's own lines, in the article's own notation, and no retained line is substituted or re-flowed.  No further source-local context is needed: every cross-reference of every retained line resolves inside the two delivered spans.  Nothing was omitted from any retained span, so the package carries no omission record.  No retained line contains a citation, so the wrapper carries no bibliography and no bibliography entry.  No retained line contains a figure environment, an image-inclusion command or drawing code, so no image is delivered and the package declares no asset dependency.  Editorial formatting only, in addition: an AMS \texttt{amsart} preamble that restates the article's own declarations and macros used by the retained lines, namely the environments \texttt{lemma} on separate counters, together with the standard packages those lines need.  The retained environments print in this document as Lemma 1 in order of appearance, whereas the article prints the same items with the numbering of its own full document.  The complete native source of the article as distributed in the arXiv e-print for this exact version is bundled unchanged as \texttt{sources/208640/source0.tex}.
\begin{lemma}\label{lma:commutatorasymptotics}
Let $p, b \in \mathbb R^\times$, let $Z = \left(\begin{smallmatrix}\alpha&\beta\\\zeta&\delta\end{smallmatrix}\right)$ be an arbitrary real $2\times 2$ matrix, and let $Y = \left(\begin{smallmatrix}\upsilon&\eta\\0&-\upsilon\end{smallmatrix}\right)$ be traceless with $Y_{21} = 0$. Put $P = \left(\begin{smallmatrix}1&p\\0&1\end{smallmatrix}\right)$ and $P_b = \left(\begin{smallmatrix}1&b\\0&1\end{smallmatrix}\right)$, and let
\begin{equation*}
	M(t) = \bigl(I + tY + t^2Y^{(2)} + O(t^3)\bigr)P \ \ \ \ \
	K(t) = \bigl(I + tZ + t^2Z^{(2)} + O(t^3)\bigr)P_b
\end{equation*}
be real analytic, with $Y^{(2)}, Z^{(2)}$ arbitrary. Then
\begin{equation*}
	M^{-1}K^{-1}MK = I + t\begin{pmatrix}p\zeta & * \\ 0 & -p\zeta\end{pmatrix} + O(t^2) \ \ \ \ \
	MK^{-1}M^{-1}K = I + t\begin{pmatrix}-p\zeta & * \\ 0 & p\zeta\end{pmatrix} + O(t^2)
\end{equation*}
Moreover,
\begin{equation*}
	\bigl(M^{-1}K^{-1}MK\bigr)_{21} = -\zeta\bigl(p\zeta + 2\upsilon\bigr)t^2 + O(t^3)
\end{equation*}
and this coefficient is independent of $\eta$, of $\alpha, \beta, \delta$, of $b$, and of $Y^{(2)}, Z^{(2)}$.
\end{lemma}

\begin{proof}
For a matrix $A(t) = A_0 + tA_1 + t^2A_2 + O(t^3)$ with $A_0$ invertible,
\begin{equation*}
	A(t)^{-1} = A_0^{-1} - tA_0^{-1}A_1A_0^{-1} + t^2\bigl(A_0^{-1}A_1A_0^{-1}A_1A_0^{-1} - A_0^{-1}A_2A_0^{-1}\bigr) + O(t^3)
\end{equation*}
Applying this to $M(t) = P + tYP + t^2Y^{(2)}P$ and simplifying with $P^{-1}(YP)P^{-1} = P^{-1}Y$ gives
\begin{equation*}
	M(t)^{-1} = P^{-1} - tP^{-1}Y + t^2\bigl(P^{-1}Y^2 - P^{-1}Y^{(2)}\bigr) + O(t^3)
\end{equation*}
and likewise for $K(t)^{-1}$ with $Y, p, Y^{(2)}$ replaced by $Z, b, Z^{(2)}$. Since $P$ and $P_b$ are both upper unitriangular they commute, so both $M^{-1}K^{-1}MK$ and $MK^{-1}M^{-1}K$ have constant term $I$ and hence begin at order $t$.
	
\medskip
	
A first-order multiplication gives
\begin{equation*}
	[t]\bigl(M^{-1}K^{-1}MK\bigr) = \begin{pmatrix}p\zeta & p(\delta-\alpha)+p^2\zeta+2b(p\zeta+\upsilon)\\ 0&-p\zeta \end{pmatrix}
\end{equation*}
and
\begin{equation*}
	[t]\bigl(MK^{-1}M^{-1}K\bigr) = \begin{pmatrix} -p\zeta & p(\alpha-\delta)+p^2\zeta-2b(p\zeta+\upsilon)\\ 0&p\zeta \end{pmatrix}
\end{equation*}
This proves the two first-order assertions.

\medskip

For the second-order calculation, write
\begin{equation*}
	M^{-1}=A_0+tA_1+t^2A_2+O(t^3) \ \ \ \ \
	K^{-1}=B_0+tB_1+t^2B_2+O(t^3)
\end{equation*}
\begin{equation*}
	M=C_0+tC_1+t^2C_2+O(t^3) \ \ \ \ \
	K=D_0+tD_1+t^2D_2+O(t^3)
\end{equation*}
using the coefficients displayed above.  Put $y_{21}^{(2)}=(Y^{(2)})_{21}$, $z_{21}^{(2)}=(Z^{(2)})_{21}$. The ten contributions to $[t^2](M^{-1}K^{-1}MK)_{21}$, listed in the order
\begin{equation*}
	\begin{gathered}
		A_2B_0C_0D_0 \ \ \ \ \
		A_0B_2C_0D_0 \ \ \ \ \
		A_0B_0C_2D_0 \ \ \ \ \ 
		A_0B_0C_0D_2 \ \ \ \ \ 
		A_1B_1C_0D_0 \\
		A_1B_0C_1D_0 \ \ \ \ \ 
		A_1B_0C_0D_1 \ \ \ \ \
		A_0B_1C_1D_0 \ \ \ \ \
		A_0B_1C_0D_1 \ \ \ \ \
		A_0B_0C_1D_1
	\end{gathered}
\end{equation*}
have $(2,1)$-entries
\begin{equation*}
	\begin{gathered}
		-y_{21}^{(2)} \ \ \ \ \
		(\alpha+\delta)\zeta-z_{21}^{(2)} \ \ \ \ \
		y_{21}^{(2)} \ \ \ \ \
		z_{21}^{(2)} \ \ \ \ \
		-\upsilon\zeta \\
		0 \ \ \ \ \
		\upsilon\zeta \ \ \ \ \
		-\upsilon\zeta \ \ \ \ \
		-(\alpha+\delta)\zeta-p\zeta^2 \ \ \ \ \
		-\upsilon\zeta
	\end{gathered}
\end{equation*}
Summing gives
\begin{equation*}
	[t^2]\bigl(M^{-1}K^{-1}MK\bigr)_{21} = -p\zeta^2-2\upsilon\zeta = -\zeta(p\zeta+2\upsilon)
\end{equation*}
This also makes explicit the cancellation of all dependence on $\eta,\alpha,\beta,\delta,b,Y^{(2)}$, and $Z^{(2)}$.
\end{proof}

\end{document}
